\section{Relative Velocity As Quotient Selection}

Relative velocity begins when two records must be compared and neither record
is allowed to own the comparison by itself. Each record carries its own local
successor count, its own order of refinements, its own boundary commitments,
and its own unresolved residue. If the grammar placed an absolute motion
behind the records before comparing them, it would undo the work of the
previous chapters. It would give the reader a view from nowhere, a place
where motion could be named without being selected by an admissible
observation.

The grammar forbids that shortcut. It permits a comparison only after the
records have been brought under a rule of observational equivalence. Two
presentations are observationally equivalent when every distinction available
to the selected reader carries the same record through both presentations.
The order of inscription may differ. The labels may shift. The local clocks
may have counted along different paths. But if the admissible comparison can
recover no difference, the presentations are identified. They fall into one
equivalence class.

Relative velocity is the residue left by that quotient. It is not the speed
that an object possesses before a reader arrives. It is the selected
difference between records after the grammar has removed the differences that
this observation cannot spend. In schematic form one may write
\[
  v_{\mathrm{rel}} = [R_1 - R_2]_{\sim},
\]
but the formula is only a reminder. The bracket is the grammar. It says that
the raw difference between records is not yet a measurement. The raw
difference must be passed through the observational quotient. What survives
is the relative velocity.

This makes the word relative stricter than ordinary perspectival talk. A
relative reading is not whatever a reader happens to prefer. The reader must
state the equivalence relation by which presentations are identified and the
selection by which one comparison is made current. Once that selection is in
place, the residue is not optional. If the quotient removes the difference,
the grammar forbids a velocity reading. If the quotient leaves a residual
count, the grammar reads relative motion.

The whole-and-parts language from the earlier count returns here. A merged
record may carry a whole comparison: the two histories are read together as
one reconciled ledger. It may also expose parts: local successor counts that
cannot be made the same without leaving a debt. Relative velocity is the
difference between the whole reconciliation and the parted counts. It reads
how much of the two local ledgers refuses to disappear when the observation
has identified everything it is allowed to identify.

This is why simultaneity cannot be primitive. Two separated readers may each
append events as local successors. A distant event that occupies one ordinal
place in the first record need not occupy the same ordinal place in the
second. A signal exchanged between them is itself an event in both ledgers,
not a magical inspection of a completed outside time. The grammar preserves
causal order where order can be transported. It does not preserve a universal
partition of distant events unless a comparison has earned that partition.
The relativity of simultaneity is the refusal to identify what no selected
observation has quotiented.

The velocity illustration says the same thing in a simpler register. A
moving mark does not carry motion as a private substance. It leaves a
sequence of distinguishable refinements in one ledger and another sequence
in a comparison ledger. When those ledgers are merged, only the difference in
recoverable counts survives. The grammar reads that surviving discrepancy as
relative velocity. If a change of labels preserves every recoverable count,
it is a presentation change. If it changes the reconciled residue, it is a
motion reading.

The Michelson--Morley-style null comparison is useful when kept in this
role. Two extremal segments may be rotated through a frame while preserving
the same refinement cost. The null result is not a declaration that a hidden
medium has failed to appear. It is a grammatical verdict: the selected
comparison cannot spend orientation as a difference in the interval count.
The quotient identifies the rotated presentations for that question. What
remains invariant is not a substance behind the records, but the cost the
grammar is forced to carry through every admissible presentation.

The velocity effect gives the opposite face. When two ledgers are reconciled
and their event counts do not identify, the discrepancy is not erased by
calling both records local. Locality tells where each record was generated;
it does not settle their comparison. The selected quotient removes label
noise and leaves the count difference. That difference is relative velocity:
motion as a carried discrepancy between records, not motion as an owned
interior property.

The selection is therefore the first act of this chapter. It permits one
comparison to stand. It forbids the reader from treating every presentation
difference as physical. It identifies observationally equivalent records. It
quotients labels, local inscriptions, and harmless rotations. It carries the
residue left after those identifications. That residue is relative motion.

Once relative velocity has been read this way, the charge problem sharpens.
A sign is also relative to an oriented comparison. The flat baseline and the
tilted baseline can read the same carried activity with opposite signs
because the selected comparison has changed. But a sign is not available on
an open local segment merely because local transport occurred. It must close
an orientation. The next section reads the positron as that closed-loop
holonomy sign.
